
\section{Feature Engineering}
\label{sec:feature-engineering}

This section describes how the raw on-chain events are transformed into the per-whale, per-day feature set used in the machine learning analysis. The pipeline
is implemented separately for Morpho Blue and Aave~v3 because the two
protocols expose different position and rate semantics, but both follow
the same logical sequence: replay events into balances, value those
balances in USD, derive risk and behavioural metrics, and finally attach
yield and rate series to compute the carry spread. The unit of
observation throughout is the \emph{whale--day}: one row per selected
whale per calendar day over which it holds an open or recently closed
position.

\subsection{Panel Construction}
\label{sec:panel-construction}

The first three stages convert raw events into a valued daily panel.
The \texttt{collect\_events} stage reads the Dune CSV exports, filters them to the supply, withdraw,
borrow, repay, and liquidation events belonging to each tracked whale,
and normalises them into a single clean event table
(\texttt{panel\_raw\_events.csv}). The \texttt{build\_panel} stage then
replays this event stream chronologically, cumulating supply and
withdraw events into a running collateral balance and borrow and repay
events into a running debt balance, emitting one row per whale per day
(\texttt{daily\_panel.csv}). Finally, \texttt{build\_prices} fetches
historical daily token prices from DeFiLlama and joins them onto the
panel, converting the native-unit collateral and debt balances into USD
values (\texttt{daily\_panel\_usd.csv}). Pricing both legs in USD is
necessary because collateral and debt are denominated in different
tokens, and risk metrics such as loan-to-value are only meaningful once
both legs share a common unit.

\subsection{Risk Features}
\label{sec:risk-features}

The risk features characterise how aggressively each position is geared
and how close it sits to liquidation on a given day. They are computed
by \texttt{build\_features} from the valued panel.

\paragraph{Loan-to-value.}
The loan-to-value ratio is the central risk measure,
\[
\mathrm{LTV} = \frac{\text{borrow}_{\text{USD}}}{\text{collateral}_{\text{USD}}}.
\]

\paragraph{Health factor.}
The health factor expresses the same relationship relative to the
market's liquidation threshold $\ell$,
$\mathrm{HF} = (\text{collateral}_\text{USD} \cdot \ell) / \text{borrow}_\text{USD}$.
A health factor below $1.0$ indicates that the position is eligible for
liquidation. The liquidation threshold is fixed per market from the
protocol's on-chain configuration
(Table~\ref{tab:liquidation-thresholds}).

\paragraph{Distance to liquidation.}
The distance to liquidation, $\ell - \mathrm{LTV}$, measures how many
percentage points of LTV buffer remain before the position becomes
liquidatable. A positive value indicates a safe position and a negative
value indicates an underwater one. Compared to the health factor, this
linear measure is more directly interpretable as a behavioural target,
since operators manage their positions in LTV space.

\paragraph{Loop depth.}
The loop depth approximates the effective leverage multiple of the
position, $\mathrm{round}\big(1 / (1 - \mathrm{LTV})\big)$, capped at the
maximum attainable for the market,
$1 / (1 - \ell)$, and set to $0$ when there is no open debt. Recursive
looping (supply, borrow, re-supply) stacks leverage that converges to
$1/(1-\mathrm{LTV})$, so this quantity reflects roughly how many times
the position is geared; an LTV of $90\%$, for example, corresponds to
about $10\times$ leverage. We derive loop depth from LTV rather than
counting discrete borrow transactions because a transaction count
overstates leverage for whales that perform frequent small rebalances
rather than discrete recursive loops---W-3, for instance, executed 170
separate borrow events while maintaining a single geared position. The
LTV-derived figure instead captures the position's effective leverage
regardless of how it was assembled. We validated this construction
against W-1, whose LTV-derived depth matched the number of
supply$\rightarrow$borrow loop transactions observed on its main build
day.

\subsection{Behavioural Features}
\label{sec:behavioural-features}

The behavioural features capture how actively each operator manages its
position over time, providing the temporal context required for the
rebalancing prediction task in RQ3.

\paragraph{Rebalancing.}
A binary flag marks each day on which any supply, withdraw, borrow, or
repay event occurred for the whale, and is zero otherwise. This flag is
the basic event indicator from which the frequency and recency features
are derived.

\paragraph{Rebalancing frequency and recency.}
Rolling counts of rebalancing events over trailing 7-day and 30-day
windows (\texttt{rebal\_freq\_7d}, \texttt{rebal\_freq\_30d}) summarise
how busy an operator has been, while \texttt{days\_since\_rebalance}
records the number of days since the last rebalancing event within the
open position. Together these distinguish high-frequency managers from
operators that set a position and leave it largely untouched.

\paragraph{LTV volatility and lags.}
The rolling 7-day and 14-day standard deviation of LTV
(\texttt{ltv\_vol\_7d}, \texttt{ltv\_vol\_14d}) captures how much the
position's gearing fluctuates, which reflects both market movement and
active management. To give the models short-horizon temporal context for
predicting rebalancing, the LTV is also lagged by 1, 3, and 7 days
(\texttt{ltv\_lag\_1}, \texttt{ltv\_lag\_3}, \texttt{ltv\_lag\_7}).

\paragraph{Token identity.}
The collateral and borrow tokens are one-hot encoded
(\texttt{coll\_*}, \texttt{debt\_*}), allowing the models to condition on
market identity, since liquidation thresholds and yield characteristics
differ by token pair.

\subsection{Yield and Carry Features}
\label{sec:carry-features}

The profitability of a looping strategy depends on the carry spread
between the yield earned on collateral and the cost of borrowing.
Computing this spread requires three rate series, joined onto the
feature table per whale by market and forward-filled within each
whale--market group to cover days without a fresh observation. On
Morpho, this is handled by a dedicated supply-rate stage
(\texttt{get\_morpho\_supply\_rates}) followed by \texttt{build\_rates};
on Aave, \texttt{build\_rates} performs the join directly.

\paragraph{Yield-bearing token APY.}
The annualised yield accruing to the collateral token (\texttt{ybt\_apy})
is sourced from the corresponding DeFiLlama pool for sUSDe, sUSDS, and
syrupUSDC.

\paragraph{Supply rate.}
The annualised protocol supply rate earned on the supplied collateral
(\texttt{supply\_rate}) is taken per market. On Aave~v3 the supply rate is 0 for the tokens we have chosen. On Morpho Blue
it is retrieved from the Morpho GraphQL API, using the
\texttt{marketById} query's \texttt{historicalState.supplyApy}
timeseries at daily resolution from a fixed historical anchor, then
deduplicated to one observation per market and day (retaining the latest
data point where a day reports more than one). We do not hardcode the
Morpho supply rate to zero: Morpho Blue markets pass through a genuine
supply APY - the native rate plus the collateral token's own yield
contribution - which we confirmed against the rates panel of the Morpho
application.

\paragraph{Borrow rate.}
The annualised borrow rate paid on the debt token (\texttt{borrow\_rate})
is taken per market: from Dune exports for the Morpho markets, and as the
per-token variable borrow rate (USDT for W-6, USDe for W-7) for the Aave
positions.

\paragraph{Carry spread.}
The per-loop carry spread combines the three rates,
$\text{carry\_spread} = \text{ybt\_apy} + \text{supply\_rate} - \text{borrow\_rate}$.
A positive spread indicates that each unit of looped exposure earns more
than it costs to fund; a negative spread indicates that the position is
losing carry and should be unwound. To express the spread at the
position's actual gearing, the levered carry spread scales it by loop
depth,
$\text{carry\_spread\_levered} = \text{carry\_spread} \cdot \text{loop\_depth}$.

\subsection{Liquidation Thresholds}
\label{sec:liquidation-thresholds}

Table~\ref{tab:liquidation-thresholds} lists the liquidation thresholds
used in the health-factor and loop-depth calculations, taken from each
market's on-chain configuration. The Morpho markets use isolated
per-pair thresholds, while the two Aave positions share Aave's sUSDe
stablecoin E-Mode threshold.

\begin{table}[t]
\centering
\caption{Liquidation thresholds by market.}
\label{tab:liquidation-thresholds}
\begin{tabular}{lll}
\hline
Protocol & Market & Threshold \\
\hline
Morpho Blue & sUSDe / USDtb & 0.915 \\
Morpho Blue & sUSDS / USDT & 0.965 \\
Morpho Blue & sUSDe / PYUSD & 0.915 \\
Morpho Blue & syrupUSDC / PYUSD & 0.915 \\
Aave v3 & sUSDe stablecoin E-Mode & 0.920 \\
\hline
\end{tabular}
\end{table}

\subsection{Feature Summary}
\label{sec:feature-summary}

Table~\ref{tab:features} summarises the full feature set produced by the
pipeline. The resulting \texttt{daily\_features} table forms the input to
the machine learning analysis, where the
risk and behavioural features characterise strategy archetypes (RQ1,
RQ2) and the lagged features support short-horizon
prediction of rebalancing and stressed risk states (RQ3).

\begin{table}[t]
\centering
\caption{Engineered features.}
\label{tab:features}
\begin{tabular}{ll}
\hline
Feature & Description \\
\hline
\texttt{ltv} & Borrow USD / collateral USD \\
\texttt{health\_factor} & $(\text{collateral USD} \times \ell)/\text{borrow USD}$; $<1.0$ liquidatable \\
\texttt{distance\_to\_liq} & $\ell - \mathrm{LTV}$; positive safe, negative underwater \\
\texttt{rebalancing} & 1 if any event occurred that day, else 0 \\
\texttt{loop\_depth} & $\mathrm{round}(1/(1-\mathrm{LTV}))$, capped at $1/(1-\ell)$; 0 if no debt \\
\texttt{ltv\_vol\_7d}, \texttt{ltv\_vol\_14d} & Rolling 7- and 14-day std.\ dev.\ of LTV \\
\texttt{rebal\_freq\_7d}, \texttt{rebal\_freq\_30d} & Rolling 7- and 30-day count of rebalancing events \\
\texttt{days\_since\_rebalance} & Days since last rebalance within the open position \\
\texttt{ltv\_lag\_1/3/7} & LTV lagged by 1, 3, and 7 days \\
\texttt{coll\_*}, \texttt{debt\_*} & One-hot identity of collateral and borrow token \\
\texttt{ybt\_apy} & Annualised yield on the collateral token (DeFiLlama) \\
\texttt{supply\_rate} & Annualised protocol supply rate on the collateral \\
\texttt{borrow\_rate} & Annualised borrow rate on the debt token \\
\texttt{carry\_spread} & $\text{ybt\_apy} + \text{supply\_rate} - \text{borrow\_rate}$ \\
\texttt{carry\_spread\_levered} & $\text{carry\_spread} \times \text{loop\_depth}$ \\
\hline
\end{tabular}
\end{table}